\documentclass[11pt,a4paper]{article}

% ---------------------------------------------------------
% PREAMBLE (stable, warning-free, Overleaf-safe)
% ---------------------------------------------------------
\usepackage[utf8]{inputenc}
\usepackage[T1]{fontenc}
\usepackage{lmodern}
\usepackage{amsmath, amssymb, amsfonts}
\usepackage{bm}
\usepackage{geometry}
\usepackage{graphicx}
\usepackage{hyperref}
\usepackage{cite}
\usepackage{authblk}
\usepackage{physics}
\usepackage{mathtools}

\geometry{margin=2.4cm}

\title{\textbf{Companion LVI: Galaxy--Galaxy Lensing in the Relaxation-Driven Cyclic Cosmology}}
\author{Michael Lehmann \\ \texttt{mi.lehmann@gmx.de}}
\date{April 2026}

\begin{document}
\maketitle

\begin{abstract}
Galaxy--galaxy lensing (GGL) probes the relationship between matter and the 
gravitational potential by measuring the distortion of background galaxies around 
foreground lenses. In the standard cosmological model, the GGL signal is determined 
by the projected matter density profile and the growth of structure. In the 
Relaxation-Driven Cyclic Cosmology (RDCC), the scale-dependent evolution of the 
gravitational potential modifies the matter distribution, the potential depth, and 
the tangential shear profile in a characteristic way. This Companion develops a 
continuous description of galaxy--galaxy lensing in RDCC, deriving how the 
relaxation scale influences the excess surface density, the tangential shear, and 
the observational signatures in weak-lensing surveys. The resulting framework 
provides a direct probe of the RDCC gravitational field on halo and large-scale 
structure scales.
\end{abstract}
% ---------------------------------------------------------
\section{Introduction}

Galaxy--galaxy lensing measures the distortion of background galaxies due to the 
gravitational potential of foreground galaxies or halos. The observable is the 
tangential shear, which is directly related to the excess surface density (ESD). 
GGL provides a powerful probe of the matter distribution, halo profiles, and the 
growth of structure.

In the standard cosmological model, the GGL signal is determined by the projected 
matter density profile and the linear growth rate. In RDCC, the gravitational 
potential evolves differently. The relaxation mechanism suppresses the decay of 
small-scale modes and enhances the coherence of large-scale modes. This modifies 
the matter distribution, the potential depth, and the tangential shear profile in a 
scale-dependent manner, producing a distinctive signature in galaxy--galaxy 
lensing.
% ---------------------------------------------------------
\section{The RDCC Matter Density Profile}

The matter density contrast evolves as
\begin{equation}
    \delta_{m,\mathrm{RDCC}}(k,z)
    = \delta_{m}(k,z)
      \left[
        1 - \chi_{\delta}
        \left(\frac{k}{k_{\mathrm{rel}}}\right)^{\alpha}
      \right].
\end{equation}

The projected surface density becomes
\begin{equation}
    \Sigma_{\mathrm{RDCC}}(R)
    = \Sigma(R)
      \left[
        1 - \chi_{\Sigma}
        \left(\frac{R_{\mathrm{rel}}}{R}\right)^{\alpha}
      \right].
\end{equation}

This modification suppresses small-scale density and enhances large-scale 
coherence.
% ---------------------------------------------------------
\section{Excess Surface Density in RDCC}

The excess surface density is defined as
\begin{equation}
    \Delta\Sigma(R)
    = \bar{\Sigma}(<R) - \Sigma(R).
\end{equation}

In RDCC, this becomes
\begin{equation}
    \Delta\Sigma_{\mathrm{RDCC}}(R)
    = \Delta\Sigma(R)
      \left[
        1 - \chi_{\Delta\Sigma}
        \left(\frac{R_{\mathrm{rel}}}{R}\right)^{\alpha}
      \right].
\end{equation}

This modification reflects the scale-dependent evolution of the gravitational 
potential.
% ---------------------------------------------------------
\section{Tangential Shear in RDCC}

The tangential shear is given by
\begin{equation}
    \gamma_{t}(R)
    = \frac{\Delta\Sigma(R)}{\Sigma_{\mathrm{crit}}}.
\end{equation}

In RDCC, this becomes
\begin{equation}
    \gamma_{t,\mathrm{RDCC}}(R)
    = \gamma_{t}(R)
      \left[
        1 - \chi_{\gamma}
        \left(\frac{R_{\mathrm{rel}}}{R}\right)^{\alpha}
      \right].
\end{equation}

This modification suppresses small-scale shear and enhances large-scale coherence.
% ---------------------------------------------------------
\section{Large-Scale Coherence and RDCC Signatures}

The relaxation mechanism enhances the coherence of large-scale modes. The GGL 
signal therefore exhibits:

\begin{itemize}
    \item enhanced shear at large radii,
    \item suppressed shear at small radii,
    \item a characteristic turnover near $R_{\mathrm{rel}}$,
    \item a modified redshift dependence of the ESD amplitude.
\end{itemize}

These signatures provide a direct observational probe of the RDCC gravitational 
field.
% ---------------------------------------------------------
\section{Conclusion}

Galaxy--galaxy lensing in the Relaxation-Driven Cyclic Cosmology reflects the 
scale-dependent evolution of the gravitational potential and the matter 
distribution. The relaxation mechanism modifies the matter density profile, the 
excess surface density, and the tangential shear in a scale-dependent manner, 
producing distinctive signatures in weak-lensing surveys. These effects provide a 
direct observational window into the relaxation scale and the temporal structure of 
the RDCC gravitational field. The present Companion establishes the theoretical 
foundation for future studies of halo structure, matter evolution, and the 
interplay between light and gravity in RDCC.

\bibliographystyle{apsrev4-2}
\nocite{*}
\bibliography{references}

\end{document}
